% Part: normal-modal-logic
% Chapter: axioms-systems
% Section: introduction

\documentclass[../../../include/open-logic-section]{subfiles}

\begin{document}

\olfileid{nml}{prf}{int}

\olsection{Pambuka}

Kita wis nduweni semantik kanggo basa modal dhasar adhedhasar
model modal, lan pangertèn yen !!a{formula} iku valid---bener ing
kabeh donya saka kabeh model---utawa valid tumrap sawijining kelas
model utawa bingkai---bener ing kabeh donya saka kabeh model ing kelas
iku, utawa kang adhedhasar bingkai kasebut. Logika lumrahe
nyambungake ciri semantis validitas kaya mangkono karo pangertèn
!!{derivability} saka teori bukti. Tujuwane yaiku ndéfinisikake
!!{derivability} ing sawijining sistem saengga !!a{formula} iku
!!{derivable} yen lan mung yen valid.

Sistem !!{derivation} kang paling prasaja lan paling tuwa
miturut sejarah diarani sistem !!{derivation} jinis Hilbert utawa
aksiomatik. Sistem !!{derivation} jinis Hilbert kanggo akeh logika
modal cukup gampang dibangun: minangka objek kajian metateoretis,
sistem kasebut prasaja (umpamane kanggo mbuktekake kasahihan lan
kelengkapan). Nanging, kanggo mbuktekake !!{formula}s, sistem mau
luwih angel digunakake tinimbang, upamane, sistem deduksi alami.

Ing sistem !!{derivation} jinis Hilbert, !!a{derivation}
saka !!a{formula} yaiku runtutan !!{formula}s kang kawiwitan saka
aksioma tartamtu, banjur liwat sawetara aturan inferensi, nganti
tekan !!{formula} kang dituju. Amarga kita pengin sistem
!!{derivation} cocog karo semantik, kita kudu njamin yen himpunan
formula kang !!{derivable} bener ing kabeh model (utawa bener ing
kabeh model kang kabeh aksiomane bener). Kaping pisan, kita bakal
misahake sawetara sipat logika modal kang dibutuhake supaya iki bisa
kelakon: logika modal ``normal''. Kanggo logika modal normal, mung ana
loro aturan inferensi kang perlu dianggep: modus ponens lan
nesesitasi. Minangka aksioma, kita njupuk kabeh instansi substitusi
saka tautologi lan, gumantung marang logika modal kang dianggep,
sawetara aksioma modal. Sanadyan kita mung nggatekake kelas kabeh
model, kabeh instansi substitusi saka~$\Ax{K}$\iftag{notprvDiamond}{}{ lan~$\Ax{Dual}$}
uga kudu dianggep aksioma. Iki wae wis ngasilake logika modal normal
paling cilik~$\Log K$.

\begin{defn}
Aturan \emph{modus ponens} yaiku skema inferensi
\begin{prooftree}
\AxiomC{$!A$}
\AxiomC{$!A \lif !B$}
\RightLabel{\MP}
\BinaryInfC{$!B$}
\end{prooftree}
Kita ngarani !!a{formula}~$!B$ \emph{ngetutake saka}~!!{formula}s
$!A$, $!C$ miturut modus ponens yen lan mung yen
$!C \ident !A \lif !B$.
\end{defn}

\begin{defn}
Aturan \emph{nesesitasi} yaiku skema inferensi
\begin{prooftree}
\AxiomC{$!A$}
\RightLabel{\Nec}
\UnaryInfC{$\Box !A$}
\end{prooftree}
Kita ngarani !!{formula}~$!B$ ngetutake saka !!{formula}s $!A$ miturut
nesesitasi yen lan mung yen $!B \ident \Box !A$.
\end{defn}

\begin{defn}
Sawijining \emph{!!{derivation}} saka himpunan aksioma~$\Sigma$
yaiku runtutan !!{formula}s $!B_1$, $!B_2$, \dots, $!B_n$, kang saben
$!B_i$ iku
\begin{enumerate}
\item instansi substitusi saka tautologi, utawa
\item instansi substitusi saka !!a{formula} ing~$\Sigma$, utawa
\item ngetutake saka loro !!{formula}s $!B_j$, $!B_k$ kanthi $j$,
  $k < i$ miturut modus ponens, utawa
\item ngetutake saka !!a{formula}~$!B_j$ kanthi $j < i$ miturut
  nesesitasi.
\end{enumerate}
Yen ana !!a{derivation} kaya mangkono kanthi $!B_n \ident !A$,
kita ngarani $!A$ \emph{!!{derivable} saka $\Sigma$}, kanthi simbol
$\Sigma \Proves !A$.
\end{defn}

Kanthi definisi iki, bakal kabukten yen himpunan formula
kang !!{derivable} mbentuk logika modal normal, lan saben
kang !!{derivable} minangka !!{formula} bener ing saben model kang kabeh
aksiomane bener. Sipat !!{derivation}s iki diarani \emph{kasahihan}.
Kosok baline, yaiku \emph{kelengkapan}, luwih angel dibuktekake.

\end{document}
